\section{Literature Review}

The starting point for the project is Taylor's classical analysis of conducting drops in electric fields \cite{taylor1964disintegration}. In the ideal electrostatic limit, the exterior potential satisfies Laplace's equation and the conical conducting boundary imposes a Legendre-function condition. Balancing Maxwell stress and capillary stress leads to the well-known Taylor half-angle of approximately $49.3^\circ$. This result is a benchmark rather than a novelty of the present work.

Subsequent studies and reviews show that real electrospray operation extends beyond the static Taylor-cone equilibrium. Fern\'andez de la Mora reviewed the fluid dynamics of Taylor cones and cone-jets, emphasizing that cone formation, emitted current, and downstream breakup are separate physical regimes \cite{fernandez2007fluid}. Melcher and Taylor developed broader electrohydrodynamic stress-balance concepts at fluid interfaces \cite{melcher1969electrohydrodynamics}, while Saville reviewed the Taylor--Melcher leaky-dielectric model, which includes finite conductivity, surface charge, and tangential electric stresses \cite{saville1997electrohydrodynamics}. Gañán-Calvo connected Taylor-type electrostatic solutions to cone-jet scaling laws \cite{ganan1997conejet}, and Fern\'andez de la Mora and Loscertales studied current emitted by highly conducting Taylor cones \cite{fernandez1994current}. Hartman and coauthors modeled cone-jet electrospray behavior with additional physical complexity \cite{hartman1999electrohydrodynamic}.

These studies motivate two important boundaries for the present project. First, the solver should reproduce the classical Taylor limit in the appropriate charge-free conducting case. Second, the solver should not overclaim full electrospray prediction. The present model does not solve transient Navier--Stokes flow, full Taylor--Melcher surface charge conservation, droplet breakup, or plume physics. Instead, it implements a reduced-order axisymmetric electrostatic--capillary model with effective Poisson space-charge closures. The intended contribution is accessibility and transparency: a sparse Python implementation that can be inspected, tested, and run on consumer hardware for onset-level parameter studies.
